% ===================================================================== Chapter 6
\chapter{CFD Methodology}\label{ch:cfdmethod}

\section{Fluent Setup}
The CFD model was set up and solved in ANSYS Fluent 2026 R1 \cite{ansysfluent} under the Student licence (4-way parallel). The case was built by a Python/Fluent
journal that writes each setting and reads it back before continuing; the settings of the final solution were read back once more from the
converged case (Table~\ref{tab:solver}).

%% generated by gen_tables.py - RE-ANALYSIS 2026
\begin{table}[htbp]
\centering
\small
\caption{Fluent solver and physical-model settings of the baseline solution}
\label{tab:solver}
\begin{tabular}{>{\raggedright\arraybackslash}p{0.30\linewidth}>{\raggedright\arraybackslash}p{0.60\linewidth}}
\toprule
Setting & Value \\
\midrule
Software & ANSYS Fluent 2026 R1 (Student licence, 4-way parallel) \\
Solver & pressure-based, coupled pressure--velocity, steady, absolute velocity \\
Energy equation & on (fluid and solid) \\
Turbulence model & $k$--$\omega$ SST, blended wall treatment (\texttt{correlation} $\omega$ wall treatment) \\
Gravity / radiation & off / off \\
Air & incompressible ideal gas, $\rho = p_{op}/(RT)$; $c_p$, $\mu$, $k$ piecewise-linear, 250--600 K \\
Inconel 718 & $\rho$ = 8190 kg/m$^3$; $k_s$, $c_{p,s}$ piecewise-linear, 293--673 K \\
Gradient & least-squares cell-based \\
Initialisation & standard: 300 K, $w$ = 23.5 m/s, 0 Pa gauge in the whole domain (cold start) \\
\bottomrule
\end{tabular}
\par\smallskip\parbox{\linewidth}{\footnotesize\raggedright Source: \texttt{06\_Fluent\_CFD/FLUENT\_SETUP\_NOTES.md} \S5--\S10 and \texttt{Baseline/CFD\_BASELINE\_RESULTS.md} \S2.}
\end{table}


\section{Solver Type}
A pressure-based solver with coupled pressure--velocity treatment was used. The maximum local Mach number is 0.097 at the outlet
centreline, so density-based compressible formulations are not required; the coupled algorithm gives robust convergence of the strongly
coupled velocity, temperature and density fields.

\section{Steady-State Formulation}
All boundary conditions are time-invariant and the physical question concerns the operating state, not the start-up transient; the problem
is therefore solved as steady (A-004). A steady solution is also the correct input for a linear static structural analysis.

\section{Energy Equation}
The energy equation is solved in both the fluid and the solid. Viscous dissipation is negligible (Brinkman number $1.8\times10^{-3}$) and
buoyancy is switched off (Richardson number $4.4\times10^{-5}$). The fluid--solid interface is a coupled wall: temperature and heat flux are
continuous and solved as part of the solution.

\section{k-omega SST}
The $k$--$\omega$ SST model \cite{menter1994} was chosen because it can be integrated through the viscous sublayer to the wall without wall
functions. In conjugate heat transfer the wall heat flux sets the solid temperature, which in turn sets the stress; a wall function would
model the quantity of interest instead of resolving it.

\section{Wall Treatment}
Fluent's blended SST wall treatment with the \texttt{correlation} $\omega$ formulation integrates the equations to the wall when the first cell
centre lies inside the viscous sublayer ($y^+ \lesssim 1$). The mesh was designed with a first fluid cell height of 12.2~\textmu m; the
converged solution confirms $y^+ \le$ 0.585 on the whole conjugate wall (Section~\ref{sec:yplus}).

\section{Material Models}
Air is an incompressible ideal gas with piecewise-linear $c_p$, $\mu$ and $k$ (8 points, 250--600~K). Inconel 718 is a new material with
constant density and piecewise-linear $k_s$ and $c_{p,s}$ (5 points, 293--673~K). Fluent's default solid material carried by the mesh file
was replaced explicitly. Temperature-dependent properties make the CFD consistent with the axially marched analytical model, which used the
same table.

\section{Boundary Conditions}
The boundary conditions of Table~\ref{tab:bc} were read back after assignment. A heat-path audit walked every wall and confirmed that the
flux of 8000~W/m$^2$ is applied only on the outer wall; any second non-zero flux would have been flagged as a double heat input. The inlet
turbulence intensity of 4.411\,\% follows the fully developed pipe estimate $I = 0.16\,Re^{-1/8}$ and is an assumption whose influence is
confined to the first few diameters.

\section{Initialization}
The whole domain, fluid and solid, was initialised from the inlet: 300~K, axial velocity 23.5~m/s and 0~Pa gauge pressure. This cold
start corresponds to the state before heating begins. No solution was initialised from another case.

\section{Numerical Schemes}
Gradients use the least-squares cell-based method and pressure is second-order throughout. The solution was advanced in stages
(Table~\ref{tab:stages}): first-order upwind for momentum, $k$, $\omega$ and energy while the flow and then the thermal field develop, followed by
second-order upwind for all equations. Only the second-order solution is used as a result.

%% generated by gen_tables.py - RE-ANALYSIS 2026
\begin{table}[htbp]
\centering
\footnotesize
\caption{Solution staging and discretisation of the baseline run}
\label{tab:stages}
\begin{tabular}{lllp{0.35\linewidth}}
\toprule
Stage & Iterations & Schemes & Purpose \\
\midrule
1 & 1--150 & first-order upwind, energy off & develop velocity and turbulence fields \\
2 & 151--300 & first-order upwind, energy on & bring up the thermal field \\
3 & 301--600 & second-order (pressure, momentum, $k$, $\omega$, energy) & production discretisation; convergence checked every 100 iterations \\
confirm & 601--700 & second-order & confirm that the converged state holds \\
\bottomrule
\end{tabular}
\par\smallskip\parbox{\linewidth}{\footnotesize\raggedright Source: \texttt{CFD\_BASELINE\_RESULTS.md} \S3. Wall clock 12 min 58 s on 4 cores for 700 iterations, including audits and exports.}
\end{table}


\section{Convergence Criteria}
Residuals alone were not accepted as evidence of convergence. The criteria combine residual targets with physical-monitor plateaus over the
last 200 iterations and with the global mass and energy balances (Table~\ref{tab:convcrit}); they were evaluated every 100 iterations of the
second-order stage.

%% generated by gen_tables.py - RE-ANALYSIS 2026
\begin{table}[htbp]
\centering
\footnotesize
\caption{Convergence criteria and final state of the baseline solution (iteration 700)}
\label{tab:convcrit}
\begin{tabular}{lll}
\toprule
Criterion & Target & Achieved \\
\midrule
Continuity residual & $< 10^{-4}$ & $9.0\times10^{-11}$ \\
Momentum residuals & $< 10^{-4}$ & $\le 1.6\times10^{-14}$ \\
$k$ / $\omega$ residuals & $< 10^{-4}$ & $7.2\times10^{-13}$ / $2.0\times10^{-13}$ \\
Energy residual & $< 10^{-6}$ & $1.6\times10^{-14}$ \\
$T_{out}$ drift over the last 200 iterations & $< 0.1$ K & $2.3\times10^{-6}$ K \\
$\Delta p$ relative drift & $< 0.1$\,\% & $6.9\times10^{-8}$\,\% \\
Interface heat-rate relative drift & $< 0.1$\,\% & $3.5\times10^{-7}$\,\% \\
Mass imbalance & $< 0.01$\,\% & $7.0\times10^{-13}$\,\% \\
Energy imbalance & $< 0.5$\,\% & $3.1\times10^{-11}$\,\% \\
\bottomrule
\end{tabular}
\par\smallskip\parbox{\linewidth}{\footnotesize\raggedright Source: \texttt{CFD\_BASELINE\_RESULTS.md} \S3; imbalances also \texttt{MASTER\_PROJECT\_DATA.csv} rows M047--M048 (VERIFIED).}
\end{table}


\section{Engineering Monitors}
Besides the residuals, fourteen physical monitors were recorded every iteration, among them the outlet bulk temperature, the maximum solid,
wetted-wall and outer-wall temperatures, the pressure drop, the heat rate across the interface and the outlet mass flow. The solid and fluid
temperature ranges were checked against the property tables at every iteration to exclude silent extrapolation.
